\begin{pastexam}[label={px:ust-112-ee-ds-19}]{臺聯大電機類 112 學年度 第 19 題}
（複選；3 分，每一選項單獨計分，不倒扣）

Which of the following statements about hashing tables is/are true?

(A)~Hashing tables use hash functions to map keys to indices in an array\\
(B)~Hashing tables may have multiple keys that map to the same index\\
(C)~Hashing tables have a fixed size and cannot grow or shrink\\
(D)~Hashing tables can use linked lists to store elements with the same hash value\\
(E)~All of the above is true
\end{pastexam}

\begin{pastsoln}
答案：(A)、(B)、(D)。

\begin{itemize}
\item (A)：對。 雜湊表的核心就是用雜湊函數 $h$ 把鍵 $k$ 對到陣列索引 $h(k)$（資結書第~12.1~節）。
\item (B)：對。 鍵的可能值遠多於表格大小，一定有不同的鍵對到同一個索引，這叫碰撞（collision）。
由鴿籠原理，鍵的範圍大於格數就避免不了。
\item (C)：錯。 負載因子太高時可以配置更大的表、把所有元素重新雜湊（rehashing），表格大小可以
改變；可延伸雜湊（extendible hashing）、線性雜湊等動態雜湊也是讓表格隨資料量成長（資結書
第~12.5~節）。
\item (D)：對。 這就是鏈結法（chaining）：每個索引掛一條串列，雜湊值相同的元素都放在同一條串列
上（資結書第~12.2~節）。
\item (E)：(C) 錯，所以錯。
\end{itemize}
\end{pastsoln}
